\chapter{Introdução}
\label{cap:introducao}

Este trabalho propõe a aproximação da ideia de autoconstrução dos sistemas vivos ao Aprendizado por Reforço Profundo. O ponto de partida é que um organismo não age apenas para maximizar um número externo. Sua atividade está ligada à manutenção de um corpo cuja continuidade não está garantida. A partir dessa inspiração, o trabalho investiga se um agente artificial pode aprender comportamentos de sobrevivência quando sua percepção e suas possibilidades de interação dependem de uma condição corporal precária.

\section{Motivação}

A Inteligência Artificial contemporânea tem alcançado resultados expressivos em tarefas de percepção, planejamento e tomada de decisão, incluindo jogos complexos, alinhamento de modelos de linguagem e manipulação robótica \cite{dota,rlhf,clothppo}. Parte importante desses avanços decorre do Aprendizado por Reforço, paradigma no qual um agente ajusta seu comportamento por meio da interação com um ambiente e das consequências expressas por sinais de recompensa \cite{SuttonBarto2018}. Quando esse paradigma é combinado com redes neurais profundas, torna-se possível aprender políticas a partir de observações de alta dimensionalidade, como imagens \cite{mnih-2013,mnih-2015}.

A recompensa oferece uma forma eficiente de especificar tarefas. Se o objetivo é fazer um agente coletar um item, uma solução direta consiste em atribuir um valor positivo à coleta. Entretanto, essa escolha não apenas orienta o aprendizado. Ela também informa previamente ao agente qual elemento do ambiente deve ser importante e qual comportamento deve ser repetido. O significado do item e a solução esperada são introduzidos pelo projetista por meio de um número.

Essa abordagem é adequada quando o interesse principal está em resolver uma tarefa bem definida. Contudo, ela se torna limitada quando o objetivo é investigar formas mais autônomas de inteligência. Um sistema vivo não recebe de um observador externo uma recompensa numérica por manter seu corpo. Ele precisa agir porque sua própria organização é precária e pode deixar de existir \cite{varela1991intentionality}. A possibilidade de morte e a necessidade contínua de manutenção fazem com que determinadas interações adquiram relevância para o organismo.

Este trabalho parte dessa diferença. Em vez de recompensar diretamente a coleta de um recurso, busca-se fazer com que o agente trabalhe para manter uma condição corporal artificial. Sua continuidade no ambiente depende de sua vida, que diminui ao longo do episódio e pode resultar em sua morte. Além disso, sua capacidade de perceber o mundo se degrada quando permanece muito tempo sem coletar recursos. O agente não deve buscar o \textit{medikit} apenas porque um número foi associado ao item. O recurso deve adquirir relevância porque altera sua condição e preserva suas possibilidades futuras de percepção, exploração e sobrevivência.

Pode-se questionar por que adotar esse caminho quando seria mais simples fornecer diretamente a recompensa desejada. A resposta está no tipo de problema investigado. A engenharia de recompensas procura descrever o comportamento que o agente deve aprender. A proposta deste trabalho procura controlar as condições a partir das quais um comportamento pode tornar-se relevante. Em vez de modificar o algoritmo para cada solução ou criar uma recompensa específica para cada ação desejada, modifica-se o corpo artificial e sua relação com o ambiente. A curiosidade permanece como um mecanismo geral de exploração, mas passa a ser modulada pelas consequências que a precariedade produz sobre a experiência do agente.

Essa mudança permite investigar se um mesmo mecanismo de aprendizado pode produzir comportamentos diferentes quando o corpo ou o ambiente são alterados. Organismos com constituições distintas não encontram exatamente o mesmo mundo de significados, mesmo quando ocupam um ambiente semelhante. Da mesma forma, alterações no ambiente modificam as possibilidades de ação disponíveis a um mesmo organismo. Trabalhos de \citeauthoronline{yuri-ambiente-rico} (\citeyear{yuri-ambiente-rico}; \citeyear{yuri-sem-objetivo}) mostram que comportamentos de navegação e forrageamento podem emergir de relações entre energia, percepção, corpo e ambiente sem a descrição direta de um objetivo comportamental. Esses estudos reforçam que a organização do agente e a riqueza de seu acoplamento com o ambiente influenciam os comportamentos que podem emergir.

A questão central desta dissertação é, portanto, se a curiosidade artificial pode ser orientada por meio do corpo do agente. Mais especificamente, investiga-se se uma condição corporal precária pode modificar a percepção e fazer com que a coleta de recursos se torne relevante sem receber uma recompensa extrínseca específica. A proposta não abandona o Aprendizado por Reforço nem elimina o uso de sinais numéricos. Ela utiliza uma recompensa intrínseca baseada em curiosidade, mas desloca a definição da solução da função de recompensa para a relação entre corpo, percepção e ambiente.

\section{Contextualização}

O Aprendizado por Reforço Profundo tem sido aplicado a diferentes classes de ambientes, desde jogos Atari baseados em imagens até tarefas de controle contínuo e simulação robótica \cite{mnih-2013,mnih-2015,ppo,clothppo}. Neste trabalho, utiliza-se a plataforma VizDoom, que oferece ambientes tridimensionais em primeira pessoa e permite que o agente aprenda a partir de observações visuais \cite{vizdoom}. Essa perspectiva é importante porque o agente não recebe uma descrição global do espaço. Ele precisa agir a partir da experiência visual produzida por sua posição e por sua orientação no ambiente.

A tarefa selecionada é a \textit{Health-Gathering}. Nela, o agente permanece em uma sala cujo chão reduz continuamente sua vida. Os \textit{medikits} distribuídos pelo ambiente restauram essa condição, e o episódio termina quando a vida se esgota. A morte do agente depende, portanto, do estado de seu corpo artificial e não de uma simples redução da recompensa acumulada. Essa dinâmica fornece uma base adequada para investigar comportamento de manutenção, pois o agente precisa encontrar recursos para continuar existindo no ambiente.

A interpretação dessa tarefa é inspirada pela Inteligência Artificial Enativa, que procura compreender a cognição a partir da organização autônoma do agente e de seu acoplamento contínuo com o ambiente \cite{FROESE2009}. Nessa perspectiva, a inteligência não consiste apenas na resolução de problemas definidos externamente. Ela emerge da atividade de um sistema corporificado que precisa manter sua própria continuidade e que produz um domínio de interações significativas a partir dessa necessidade.

O conceito de autopoiese descreve os sistemas vivos como redes de processos que produzem e mantêm continuamente os próprios componentes \cite{maturana1980autopoiesis,varela1991intentionality}. A organização autopoiética não está garantida de antemão. Ela permanece exposta à possibilidade de desintegração e precisa ser reconstruída por meio da própria atividade do organismo. Essa condição é denominada precariedade. Um sistema precário depende de processos que precisam continuar ocorrendo para que sua identidade seja preservada.

O agente desenvolvido nesta dissertação não realiza autopoiese em sentido próprio. Ele não produz seus componentes, não define sua identidade e não estabelece autonomamente seus limites de viabilidade. O glaucoma artificial oferece apenas uma aproximação operacional à ideia de um corpo que precisa ser mantido. Quando o agente permanece sem coletar \textit{medikits}, sua percepção se degrada progressivamente. Quando coleta um recurso, recupera a condição perceptiva. Assim, agir sobre o ambiente produz uma forma limitada de reconstrução de sua capacidade de continuar percebendo e explorando.

Essa relação transforma o \textit{medikit}. Do ponto de vista de um observador externo, ele continua sendo um objeto presente no cenário. Para o agente, entretanto, o item pode adquirir significado funcional porque interfere em sua condição corporal e em suas experiências futuras. A precariedade faz com que o ambiente deixe de ser apenas um conjunto de elementos visuais e passe a constituir um mundo de interações com consequências diferentes para a continuidade do agente.

O mecanismo utilizado para explorar essa relação é o glaucoma artificial. Uma região escurecida cresce progressivamente sobre a observação enquanto o agente permanece sem coletar recursos. A velocidade desse crescimento é controlada pela força do glaucoma. Quando um \textit{medikit} é coletado, a observação é restaurada. O glaucoma não busca reproduzir clinicamente a doença humana. Ele funciona como uma forma computacional de precariedade perceptiva que conecta o estado do corpo artificial à capacidade de acessar o ambiente.

A política é treinada com o algoritmo \textit{Proximal Policy Optimization} (PPO) \cite{ppo}. A curiosidade é implementada por meio do \textit{Random Network Distillation} (RND), que utiliza o erro de uma rede preditora como medida de novidade das observações \cite{rnd}. Observações frequentes e degradadas tendem a tornar-se previsíveis e produzir recompensas menores. Quando a visão é restaurada, a imagem volta a apresentar maior quantidade de informação visual e pode produzir maior erro de predição. Dessa maneira, o corpo artificial modula a curiosidade sem que seja necessário atribuir uma recompensa específica ao \textit{medikit}, à restauração da visão ou à sobrevivência.

O deslocamento proposto ocorre da engenharia da solução por recompensa para a engenharia das condições corporais e ambientais. O PPO e o RND permanecem os mesmos nas diferentes forças positivas de glaucoma. O que muda é a velocidade com que a condição perceptiva se deteriora. Se as políticas aprendidas também mudarem, essa diferença não decorrerá de uma nova regra de aprendizado ou de uma nova recompensa associada à coleta. Ela decorrerá da alteração do corpo e de seu acoplamento com o ambiente.

Essa proposta constitui uma contribuição na direção da autonomia e da inteligência artificial, mas não pretende resolver o problema da autonomia. O trabalho não avalia autonomia constitutiva nem adaptatividade em seu sentido forte. Seu objetivo é mais restrito. Busca-se investigar uma forma de trazer para o Aprendizado por Reforço a ideia de que a atividade de um agente pode estar relacionada à manutenção de sua própria condição. Trata-se de um pequeno passo inspirado pela autoconstrução da vida, realizado dentro de uma arquitetura que permanece projetada externamente.

\section{Proposta}

A proposta experimental combina três condições de recompensa e quatro configurações corporais. A primeira condição não utiliza recompensa e funciona como controle negativo. A segunda utiliza a recompensa extrínseca da tarefa e funciona como controle positivo, pois informa diretamente ao agente quais consequências devem ser valorizadas. A terceira utiliza apenas a recompensa intrínseca calculada pelo RND e corresponde à proposta central do trabalho.

As forças de glaucoma avaliadas são $0$, $50$, $100$ e $200$. A força $0$ preserva a visão durante todo o episódio e funciona como uma ablação da precariedade perceptiva. As demais forças produzem degradação em velocidades crescentes. Essa comparação permite observar se o comportamento muda quando o mesmo agente, treinado com o mesmo algoritmo e com o mesmo mecanismo de curiosidade, passa a possuir condições corporais diferentes.

A hipótese investigada é que a curiosidade isolada não estabelece uma relação confiável entre os \textit{medikits} e a continuidade do agente. Quando a percepção se degrada, entretanto, a ausência de coleta torna as observações progressivamente mais pobres e previsíveis. A restauração da visão produzida pelo recurso volta a disponibilizar um ambiente visualmente rico. O RND pode então fornecer recompensas maiores em estados menos precários, favorecendo indiretamente ações que mantêm a capacidade perceptiva.

Espera-se, assim, que o \textit{medikit} adquira relevância não porque a recompensa descreve sua importância, mas porque o item transforma a relação entre o agente e o ambiente. A política continua sendo atualizada por um sinal numérico, como exige o Aprendizado por Reforço utilizado nesta pesquisa. Contudo, o vínculo entre esse sinal e a sobrevivência não é programado como uma recompensa direta pela coleta. Ele depende da interação entre corpo artificial, degradação perceptiva, novidade das observações e história de aprendizado.

Os resultados são avaliados por meio da duração média dos episódios, da variabilidade entre sementes, das trajetórias em diferentes distribuições de recursos e da relação entre recompensa intrínseca e precariedade. A análise é complementada pelo \textit{Gradient-weighted Class Activation Mapping} (Grad-CAM), que permite identificar regiões das observações relevantes para as decisões das políticas \cite{gradcam}. A combinação dessas dimensões busca verificar não apenas se os agentes sobrevivem, mas também como navegam, como a condição corporal modifica a curiosidade e quais elementos visuais participam das decisões.

\section{Objetivos}

O objetivo geral deste trabalho é investigar experimentalmente se uma forma artificial de precariedade corporal pode orientar a curiosidade e favorecer comportamentos de manutenção em agentes de Aprendizado por Reforço Profundo sem recompensar diretamente a coleta de recursos. Nesse contexto, comportamento de manutenção designa uma resposta operacional às alterações de percepção e viabilidade da tarefa. O termo não indica a implementação de autopoiese, autonomia constitutiva ou adaptatividade em sentido forte.

De forma específica, este trabalho tem como objetivos

\begin{itemize}
	\item Estruturar uma tarefa experimental na qual a continuidade do episódio dependa do estado corporal do agente e da coleta de recursos
	\item Implementar um mecanismo de glaucoma artificial capaz de degradar e restaurar gradualmente a percepção conforme a interação com os \textit{medikits}
	\item Relacionar a condição corporal artificial à recompensa por curiosidade sem atribuir uma recompensa específica à coleta, à sobrevivência ou à restauração da visão
	\item Treinar agentes sob ausência de recompensa, recompensa extrínseca e recompensa intrínseca baseada em RND
	\item Avaliar como diferentes forças de glaucoma modificam o aprendizado, a estabilidade e a capacidade de sobrevivência
	\item Analisar se alterações no corpo artificial produzem mudanças no comportamento sem exigir modificações no PPO ou no mecanismo de curiosidade
	\item Comparar as trajetórias diante de diferentes distribuições espaciais de \textit{medikits}
	\item Investigar a relação temporal entre precariedade perceptiva e recompensa intrínseca
	\item Examinar quais regiões das observações influenciam as decisões das políticas aprendidas
\end{itemize}

\section{Organização}

Além desta introdução, o trabalho está organizado em cinco capítulos. O Capítulo~\ref{cap:fundamentacao-teorica} apresenta a fundamentação filosófica e computacional da pesquisa, abrangendo inteligência, autonomia, autopoiese, precariedade, Aprendizado por Reforço Profundo, motivação intrínseca, curiosidade artificial e interpretabilidade visual. O Capítulo~\ref{cap:trabalhos-relacionados} discute estudos sobre Aprendizado por Reforço Profundo, engenharia de recompensas, aprendizado enativo e aprendizado orientado por curiosidade. O Capítulo~\ref{chap:metodologia} descreve a tarefa experimental, o corpo artificial do agente, o mecanismo de glaucoma, os algoritmos de aprendizado, as condições experimentais e as métricas de avaliação. O Capítulo~\ref{chap:resultados} apresenta e discute os resultados quantitativos e qualitativos. Por fim, o Capítulo~\ref{chap:conclusoes-e-trabalhos-futuros} reúne as conclusões da pesquisa, suas contribuições, suas limitações e as possibilidades de continuidade.
